% Generated draft: requires mathematical and source review.
\documentclass[11pt]{article}
\usepackage[margin=25mm]{geometry}
\usepackage{fontspec}
\setmainfont{Latin Modern Roman}
\usepackage{amsmath,amssymb,mathtools,unicode-math}
\setmathfont{Latin Modern Math}
\usepackage{enumitem,needspace,xurl,hyperref}
\setlength{\parindent}{0pt}
\setlength{\parskip}{5pt}
\setlength{\emergencystretch}{3em}
\newcommand{\N}{\mathbb{N}}
\newcommand{\Z}{\mathbb{Z}}
\newcommand{\Q}{\mathbb{Q}}
\newcommand{\R}{\mathbb{R}}
\newcommand{\C}{\mathbb{C}}
\newcommand{\F}{\mathbb{F}}
\newcommand{\A}{\mathbb{A}}
\newcommand{\PP}{\mathbb{P}}
\newcommand{\E}{\mathbb{E}}
\newcommand{\eps}{\varepsilon}
\newcommand{\dd}{\,\mathrm{d}}
\begin{document}
\section*{UnsolvedMath ID: 30006511}
\textbf{Source problem code:} catalogue entry
\subsection{Definitions and mathematical statement}
A closed hyperbolic surface is a compact, connected, smooth surface without boundary, equipped with a Riemannian metric of constant Gaussian curvature \(-1\). Under the usual interpretation of the source notation \(\mathcal M_g\), restrict to orientable surfaces of genus \(g\geq 2\), and let \(\mathcal M_g\) denote their isometry classes. Orientability is implicit in this notation rather than explicit in the supplied corrected statement; this interpretation requires review against the report.

Use the nonnegative Laplace--Beltrami operator \(\Delta_X=-\operatorname{div}_X\nabla_X\). Write its eigenvalues, counted with multiplicity, as
\[
0=\lambda_0(X)<\lambda_1(X)\leq\lambda_2(X)\leq\cdots.
\]
Thus \(\lambda_1(X)\) is the first positive Laplace eigenvalue.

The supplied question is whether
\[
\forall G\in\Z,\ G\geq 2,\quad
\exists g\in\Z,\ g\geq G,\quad
\exists X\in\mathcal M_g
\quad\text{such that}\quad
\lambda_1(X)>\frac14.
\]
The inequality is strict for each selected surface. No positive lower bound on \(\lambda_1(X)-\tfrac14\) independent of genus is requested. The source calls the threshold asymptotically optimal, but supplies no precise formulation of that qualification; it is retained as source commentary, not an additional asserted theorem. The registry labels the problem open; that label is supplied curation, not an independently verified mathematical conclusion.
\subsection{Short English statement}
Can one find closed orientable hyperbolic surfaces with unbounded genus whose first positive Laplace eigenvalue exceeds \(1/4\)? The curvature normalization is \(-1\), and the excess over \(1/4\) may depend on the surface. Orientability follows the usual reading of the source notation \(\mathcal M_g\) and remains subject to source review.
\subsection{Sources}
\begin{itemize}
\item Local registry: \texttt{data/unsolvedmath/problems.json}, permanent UnsolvedMath ID \(30006511\), source problem code \texttt{catalogue entry}; collection: Oberwolfach Reports Open Problems. The project also identifies \texttt{catalog/problems.json} as a registry of additional existing records. This draft uses only the supplied complete record.
\item Supplied report citation: \emph{Mini-Workshop: Hyperbolic meets Stochastic Geometry} (2026), Oberwolfach Reports, pp. 3031--3062, DOI 10.4171/owr/2025/57. \url{https://doi.org/10.4171/owr/2025/57}. The original excerpt identifies p. 3051. No individual proposer is supplied. The record separately gives proposed year 2025.
\item Mondal, \emph{On largeness and multiplicity of the first eigenvalue}. \url{https://arxiv.org/abs/1406.1080}.
\item Fortier Bourque--Petri, \emph{Linear programming bounds for hyperbolic surfaces}. \url{https://arxiv.org/abs/2302.02540}.
\item Wu--Xue, \emph{Random hyperbolic surfaces}. \url{https://arxiv.org/abs/2102.05581}.
\item Literature assessment dated exactly \texttt{2026-08-22}: the supplied curation labels the question open and states that large-genus constructions approach the one-quarter threshold from below, while existence in arbitrarily large genus above that threshold remains unresolved. This is supplied dated curation derived from the OWR record, not independently verified here; its time-sensitive status and cited claims require independent verification.
\item The registry marks the statement \texttt{corrected\_verified} and describes the repairs as changes to formula delimiters, notation, grammar, and typography preserving mathematical content and scope. This is a supplied verification claim, not independent source verification by this draft.
\end{itemize}
\end{document}
